In dimension two take $a=(0,0)$ and define $f(x,y)=x^3/(x^2+y^2)$ off zero, with $f(0,0)=0$. It is continuous everywhere because $|f(x,y)|\le |x|$, and $f(t v)=t f(v)$ for every real $t$. For any differentiable curve through zero, $\gamma(t)/t\to\gamma'(0)$. Consequently $f(\gamma(t))/t=f(\gamma(t)/t)\to f(\gamma'(0))$ from both sides, including when the velocity is zero. All the requested scalar derivatives exist. A Fréchet derivative at zero would have to take each $v$ to $f(v)$ by testing straight lines. But $f(1,0)=1$, $f(0,1)=0$ and $f(1,1)=1/2$ contradict linearity. Thus the implication fails even on the open set $\mathbb R^2$.
